%=======================================================================
%  CSE558 Data Science - Midterm Project Report (course template)
%  Monsoon 2026, IIIT-Delhi  --  Group Prompters
%  Compile with pdfLaTeX (twice) or tectonic.
%=======================================================================
\documentclass[10pt,a4paper,twocolumn]{article}

%========================= EDIT THIS =========================
\newcommand{\projecttitle}{Predictive Flight Fare Dynamics: Identifying Optimal
  Booking Horizons and Date-Specific Pricing}
\newif\ifshowhints
\showhintsfalse
%=============================================================

%-------------------------------------------- fonts (no need to edit)
\usepackage{iftex}
\ifPDFTeX
  \usepackage[T1]{fontenc}
  \usepackage[utf8]{inputenc}
  \IfFileExists{carlito.sty}%
    {\usepackage[sfdefault]{carlito}}%              Calibri look
    {\usepackage[scaled=0.92]{helvet}\renewcommand{\familydefault}{\sfdefault}}%
  \usepackage{courier}
\else
  \usepackage[no-math]{fontspec}
  \setmainfont{Carlito}[Extension=.ttf, UprightFont=*-Regular, BoldFont=*-Bold,
    ItalicFont=*-Italic, BoldItalicFont=*-BoldItalic]
\fi

%----------------------------------------------------------- packages
\usepackage[margin=18mm,bottom=20mm,columnsep=7mm]{geometry}
\usepackage{amsmath}
\usepackage{xcolor}
\usepackage{graphicx}
\usepackage{tabularx}
\usepackage{array}
\usepackage{booktabs}
\usepackage{enumitem}
\usepackage{titlesec}
\usepackage{needspace}
\usepackage[font=small,labelfont=bf]{caption}
\usepackage[font=footnotesize,labelfont=bf,skip=1pt]{subcaption}
\usepackage[hidelinks]{hyperref}
\hypersetup{pdftitle={\projecttitle},
            pdfsubject={CSE558 Data Science, Monsoon 2026: Midterm Project Report}}

%------------------------------------------------------------- layout
\definecolor{eg}{HTML}{7396C4}         % worked example (light blue)
\definecolor{hintgrey}{gray}{0.58}     % instructions (light grey)

\titleformat{\section}{\normalfont\large\bfseries}{\thesection.}{0.5em}{}
\titlespacing*{\section}{0pt}{12pt}{3pt}
\titleformat{\subsection}{\normalfont\normalsize\bfseries}{\thesubsection}{0.5em}{}
\titlespacing*{\subsection}{0pt}{10pt}{2pt}
\newcommand{\sectionbreak}{\needspace{10\baselineskip}}
\newcommand{\subsectionbreak}{\needspace{20\baselineskip}}

\setlength{\parindent}{0pt}
\setlength{\parskip}{4pt}
\setlist{leftmargin=1.2em, itemsep=2pt, topsep=2pt, parsep=0pt}
\renewcommand{\arraystretch}{1.25}
\setlength{\tabcolsep}{4pt}
\newcolumntype{L}[1]{>{\raggedright\arraybackslash}p{#1}}
\newcolumntype{Y}{>{\raggedright\arraybackslash}X}
\newcolumntype{K}[1]{>{\raggedright\arraybackslash\bfseries\boldmath}p{#1}}  % bold labels

%--------------------------------------------------------- shorthands
\newcommand{\hint}[1]{\ifshowhints{\footnotesize\itshape\color{hintgrey}#1\par}\fi}
\newcommand{\eg}[1]{#1}
\newcommand{\egblock}[1]{{#1\par}}
\newcommand{\field}[2]{\textbf{#1}~\eg{#2}\par}
\newcommand{\tpart}[2]{\par\needspace{12\baselineskip}\smallskip\textbf{#1)~#2}\par}
\newcommand{\minihead}[1]{\par\needspace{8\baselineskip}\smallskip\textbf{#1}\par}
\newcommand{\code}[1]{\texttt{#1}}
\newcommand{\figfile}[2]{%
  \IfFileExists{#1}%
    {\includegraphics[width=\linewidth,height=#2,keepaspectratio]{#1}}%
    {\fbox{\parbox[c][#2][c]{\dimexpr\linewidth-2\fboxsep-2\fboxrule\relax}%
      {\centering\small\itshape\color{hintgrey}Your plot: \detokenize{#1}}}}}
\newcommand{\rs}{Rs.\,}
% One sub-plot with its one-line takeaway
\newcommand{\subfig}[4]{\begin{subfigure}[t]{#1\linewidth}\centering
  \figfile{figures/#2}{#3}\caption{#4}\end{subfigure}}

%=======================================================================
\begin{document}

\twocolumn[%
{\small\bfseries CSE558 Data Science\ \ \textbullet\ \ Midterm Project Report\ \ \textbullet\ \ Monsoon 2026\par}
\vspace{4pt}
{\LARGE\bfseries\raggedright\projecttitle\par}
\vspace{6pt}
\rule{\linewidth}{1pt}

{\small
\begin{tabularx}{\linewidth}{@{}K{2.4cm} Y@{}}
\toprule
Group name & Prompters (Text / Documents, Tabular / Transactional, Finance / Economics) \\
Dataset & \href{https://www.kaggle.com/datasets/shubhambathwal/flight-price-prediction}{kaggle.com/datasets/shubhambathwal/flight-price-prediction} \\
Code & \href{https://github.com/Viren286/DSC_PROJECT}{github.com/Viren286/DSC\_PROJECT} (notebook \code{main\_v2.ipynb}, \code{hypothesis\_tests.py} and README) \\
\bottomrule
\end{tabularx}}

\vspace{2pt}
{\small
\begin{tabularx}{\linewidth}{@{}L{0.5cm} L{3.4cm} L{2.3cm} L{4.8cm} Y@{}}
\toprule
\textbf{\#} & \textbf{Name} & \textbf{Roll no.} & \textbf{Institute email} & \textbf{Contribution} \\
\midrule
1 & Viren Poswal & 2024623 & viren24623@iiitd.ac.in & Dataset and size check (1) \\
2 & Sahitya Sharma & 2023464 & sahitya23464@iiitd.ac.in & Preparation, before and after (2, 3) \\
3 & Aadit Garg & 2024003 & aadit24003@iiitd.ac.in & EDA and problem statements (4) \\
4 & Madhav Raj Gaur & 2024323 & madhav24323@iiitd.ac.in & Tests 1--3 (5) \\
5 & Shivam Harit & 2024528 & shivam24528@iiitd.ac.in & Tests 4--5, limitations (5, 6) \\
\bottomrule
\end{tabularx}}
\vspace{12pt}
]

%-----------------------------------------------------------------------
\section{Dataset}

\field{Source:}{Kaggle, ``Flight Price Prediction'': flight listings scraped from the
EaseMyTrip booking portal for the six largest Indian metros (Delhi, Mumbai, Bangalore, Kolkata,
Hyderabad, Chennai) over a 50-day window. One row is one flight listing on one search date.}

\field{What is raw about it:}{a redundant row index (\code{Unnamed: 0}); ordinal and nominal
attributes stored as text (\code{departure\_time}, \code{stops}, \code{class}); a right-skewed,
bimodal target (\rs1,105 to \rs123,071); and a flight-code string with 1,561 levels.}

{\small
\begin{tabularx}{\linewidth}{@{}L{4.3cm} Y Y@{}}
\toprule
 & \textbf{Raw file} & \textbf{After encoding} \\
\midrule
Instances (rows) & 300{,}153 & 297{,}920 \\
Attributes (columns) & 12 & 1{,}629 \\
Instances $\times$ attributes & 3{,}601{,}836 & 485{,}311{,}680 \\
\bottomrule
\end{tabularx}}

\field{Size check:}{$297{,}920\times1{,}629=485{,}311{,}680\ge3{,}000{,}000$, so the
requirement is met (the raw file alone has $300{,}153\times12=3{,}601{,}836$). One-hot encoding
of the 1,561 flight codes and 30 engineered routes is what takes it to 1,629 columns. The
modelling matrix after dropping \code{flight} is $297{,}920\times68=20{,}258{,}560$, still
above 3M.}

\minihead{Key features}
{\footnotesize
\begin{tabularx}{\linewidth}{@{}L{2.2cm} L{1.5cm} Y@{}}
\toprule
\textbf{Feature} & \textbf{Type} & \textbf{Real-world meaning} \\
\midrule
\code{price} & continuous & ticket fare in INR; our response \\
\code{days\_left} & discrete & days between booking and departure (1--49) \\
\code{duration} & continuous & total travel time in hours, including layovers \\
\code{class} & ordinal & Economy $<$ Business seat \\
\code{airline} & nominal & operating carrier, 6 levels \\
\code{departure\_time} & ordinal & time-of-day slot, 6 levels \\
\code{stops} & ordinal & zero $<$ one $<$ two or more \\
\bottomrule
\end{tabularx}}

%-----------------------------------------------------------------------
\section{Data preparation}

{\footnotesize
\begin{tabularx}{\linewidth}{@{}L{1.9cm} L{2.4cm} Y@{}}
\toprule
\textbf{Issue found} & \textbf{What we did} & \textbf{Why, and not the alternative} \\
\midrule
Redundant index \code{Unnamed: 0}
 & dropped
 & a running row number carries no information and would leak row order into a model \\
Missing values, duplicates
 & checked: 0 missing, 0 duplicate rows
 & nothing to impute or drop \\
2,110 \code{duration} values above the IQR fence (30.18\,h)
 & rows dropped ($Q_3+1.5\,\mathrm{IQR}$ rule)
 & long-layover itineraries inflate fare variances, which Test 3 compares. Capping would invent
   durations that were never flown \\
123 \code{price} values above the fence (\rs99,128)
 & rows dropped
 & 0.04\% of rows, all extreme Business fares, that would dominate the variance tests; a log
   transform would change the scale of every test \\
\code{flight} has 1,561 levels
 & dropped before modelling (Fig.~\ref{fig:eda}j)
 & 1,561 sparse dummy columns, and the code prefix is fully determined by \code{airline} \\
Text categoricals
 & one-hot encoded; added \code{route} = source$\to$destination
 & linear models need numbers; label codes would impose a false order on airlines and cities \\
\bottomrule
\end{tabularx}}

%-----------------------------------------------------------------------
\section{What changed after preparation}

{\footnotesize
\begin{tabularx}{\linewidth}{@{}L{2.9cm} L{1.9cm} Y@{}}
\toprule
\textbf{Measure} & \textbf{Before} & \textbf{After} \\
\midrule
Rows & 300{,}153 & 297{,}920 (2{,}233 dropped, 0.74\%) \\
Attributes & 12 & 11 raw; 1{,}629 encoded; 68 for modelling \\
Missing / duplicate rows & 0 / 0 & 0 / 0 \\
Duration range, mean & 0.83--49.83\,h, 12.22\,h & 0.83--30.17\,h, 12.07\,h \\
Price range, median & \rs1,105--123,071, \rs7,425 & \rs1,105--98,972, \rs7,424 \\
Economy share & 68.85\% & 68.74\% \\
\bottomrule
\end{tabularx}}

\egblock{Dropping 0.74\% of rows cut the longest trip from 49.8\,h to 30.2\,h and the highest
fare from \rs123k to \rs99k, while the median fare moved by \rs1 and the class mix by 0.1
percentage points. Dropping \code{flight} shrank the modelling matrix from 1,629 to 68 columns
without losing airline information.}

%-----------------------------------------------------------------------
\section{Basic questions and problem statement}

\begin{figure*}[!t]
\centering
\subfig{0.325}{fig01.png}{4.3cm}{Vistara and Air India one-stop fares have the widest, highest spread; budget carriers stay below \rs35k.}\hfill
\subfig{0.325}{fig02.png}{4.3cm}{Fares form two bands (Economy, Business) and the fitted line slopes gently down with days left.}\hfill
\subfig{0.325}{fig03.png}{4.3cm}{Every departure slot is bimodal; Late Night lacks most of the Business mode.}\\[3pt]
\subfig{0.19}{fig04.png}{4.3cm}{Numeric features are only weakly correlated: no multicollinearity.}\hfill
\subfig{0.17}{fig06.png}{4.3cm}{68.7\% of listings are Economy, 31.3\% Business.}\hfill
\subfig{0.29}{fig05.png}{4.3cm}{Vistara and Air India average 4--7$\times$ the fare of low-cost carriers.}\hfill
\subfig{0.31}{fig07.png}{4.3cm}{Mean fare is flat to about 20 days out, climbs steeply, peaks 2 days out, then drops on the last day.}\\[3pt]
\subfig{0.325}{fig08.png}{4.3cm}{Price is far from normal (QQ $R^2=0.756$), so every parametric test gets a rank-based or resampling check.}\hfill
\subfig{0.325}{fig09.png}{4.3cm}{Price density splits cleanly by class; duration is multimodal (mean 12.07\,h).}\hfill
\subfig{0.325}{fig10.png}{4.3cm}{\code{flight} has 1,561 levels and its code prefix maps one-to-one to \code{airline}, so we drop it.}
\caption{Exploratory data analysis of the cleaned data ($N=297{,}920$). Each panel has its one-line takeaway.}
\label{fig:eda}
\end{figure*}

\minihead{Basic questions}
\begin{enumerate}[label=\arabic*.]
  \item Which airline is cheapest? AirAsia (\rs4,091 on average); Vistara is dearest (\rs30,369).
  \item How far apart are Economy and Business? \rs6,530 vs.\ \rs52,446 on average, an
  8$\times$ gap that makes the fare distribution bimodal.
  \item Does a longer trip cost more? Weakly: Pearson $r=0.215$ (and $r=-0.091$ for
  \code{days\_left}).
  \item Which departure slot is dearest? Night (\rs23,037); Late Night is cheapest (\rs9,295).
  \item Which routes are busiest? Delhi$\to$Mumbai and Mumbai$\to$Delhi, with median fares of
  \rs7,262 and \rs6,300.
\end{enumerate}

\minihead{Our problem statements}
\noindent\fbox{\begin{minipage}{\dimexpr\linewidth-2\fboxsep-2\fboxrule\relax}
\textbf{P1 (primary, regression).} How does the fare of a flight change with the number of days
left before departure, and after how many days does booking earlier stop saving money (the
booking ``sweet spot'')? Does the last day bring a drop?

\smallskip
\textbf{P2.} Are late bookings ($\le20$ days) more expensive \emph{and} more volatile than early
bookings ($>20$ days), and is that explained by a different class mix?
\end{minipage}}

\field{How they are novel:}{the basic questions describe static averages of one attribute at a
time. Ours follow the fare \emph{trajectory} across the booking horizon, locate where the
advance-purchase saving stops, and separate the horizon effect from the class mix, which drives
most of the raw variance.}

%-----------------------------------------------------------------------
\section{Statistical inferences}

\field{Rows tested:}{all 297,920 cleaned listings. Late $=$ \code{days\_left} $\le20$
(112,403 rows), early $>20$ (185,517); last-minute $=\le7$ days. Every plan in part a) was fixed
before the test was run. With five tests we also check each decision at the Bonferroni level
$0.05/5=0.01$; none changes. With $N\approx3\times10^5$ almost any gap is significant, so we judge
each test by whether its 95\% CI clears the smallest difference that matters.}

{\footnotesize
\begin{tabularx}{\linewidth}{@{}L{0.25cm} Y L{1.45cm} L{1.2cm} L{1.0cm}@{}}
\toprule
\textbf{\#} & \textbf{Question} & \textbf{Statistic} & \textbf{Decision} & \textbf{Matters?} \\
\midrule
1 & Do fares drop on the last day? & $t=14.96$ & reject $H_0$ & yes \\
2 & Are late bookings dearer? & $Z=46.09$ & reject $H_0$ & yes \\
3 & Are late fares more volatile? & $F=0.977$ & fail to reject & no (pooled) \\
4 & Does fare fall with days left? & $t=-49.77$ & reject $H_0$ & yes \\
5 & Is class mix tied to timing? & $\chi^2=46.89$ & reject $H_0$ & no \\
\bottomrule
\end{tabularx}}


%------------------------------------------------------------ Test 1
\needspace{22\baselineskip}\subsection*{Test 1 (P1): do fares drop on the last day?}

\tpart{a}{Hypotheses, written before running the test}
{\small
\begin{tabularx}{\linewidth}{@{}K{3.4cm} Y@{}}
\toprule
$H_0$ & $\mu_{2}-\mu_{1}\le0$ \\
$H_1$ & $\mu_{2}-\mu_{1}>0$ \\
Significance level & $\alpha=0.05$ \\
Side & one-sided (right) \\
Test & Welch t-test \\
Smallest difference that matters & \rs1,000 \\
\bottomrule
\end{tabularx}}

\egblock{$\mu_2$, $\mu_1$: mean fare with 2 days and 1 day left. One-sided, because Fig.~\ref{fig:eda}g
suggests a last-day clearance drop. \rs1,000 is about 5\% of a typical ticket, a saving a
traveller would notice.}

\tpart{b}{Other tests we considered}
\egblock{Two independent groups (3,942 and 1,884 listings), a continuous response and unknown
$\sigma$ call for a two-sample t-test. We use Welch's version because the spreads differ
(SD \rs24,138 vs.\ \rs18,841; Levene $p=10^{-50}$). With $n>1{,}800$ per group the CLT makes the
mean difference approximately normal despite the skew.}
{\footnotesize
\begin{tabularx}{\linewidth}{@{}L{2.4cm} Y@{}}
\toprule
\textbf{Also considered} & \textbf{Why we did not use it} \\
\midrule
Pooled t-test & assumes equal variances, which Levene rejects. Run anyway: $t=13.72$ (5,824), $p=2\times10^{-42}$, same decision \\
Mann--Whitney U & tests a shift in the whole distribution, not the mean fare paid. Agrees: $p=4\times10^{-24}$, $P(\text{day 2}>\text{day 1})=0.58$, medians \rs17,854 vs.\ \rs15,078 \\
Permutation test & assumption-free check of the mean gap; 10,000 shuffles give $p<10^{-4}$ \\
Bootstrap CI & assumption-free CI for the gap: [\rs7,542, \rs9,759], matching Welch \\
\bottomrule
\end{tabularx}}

\tpart{c}{Result}
{\small
\begin{tabularx}{\linewidth}{@{}K{3.4cm} Y@{}}
\toprule
Estimate & \rs8,672 (\rs30,024 vs.\ \rs21,351) \\
95\% confidence interval & \rs7,535 to \rs9,809 \\
Test statistic (d.f.) & $t=14.96$ (4,634) \\
p-value & $1.0\times10^{-49}$ \\
Decision & reject $H_0$ \\
\bottomrule
\end{tabularx}}

\egblock{Fares drop by about \rs8.7k on the final day (Cohen's $d=0.38$). The whole interval is
more than seven times the \rs1,000 threshold, so the drop matters (Fig.~\ref{fig:tests}a).}

\begin{figure*}[!t]
\centering
\subfig{0.24}{fig11.png}{4.5cm}{Test 1: mean fare peaks 2 days out and falls sharply on the final day.}\hfill
\subfig{0.24}{fig12.png}{4.5cm}{Test 2: late bookings have a higher median and upper quartile.}\hfill
\subfig{0.24}{fig13.png}{4.5cm}{Test 4: the fitted line falls by about \rs152 per extra day booked ahead.}\hfill
\subfig{0.24}{fig15.png}{4.5cm}{Late Night fares sit lowest; Night and Morning highest.}\\[3pt]
\subfig{0.22}{fig14.png}{4.5cm}{31.7\% of tickets cost over \rs20,000, nearly all Business.}\hfill
\subfig{0.26}{fig16.png}{4.5cm}{Test 3: pooled fare SD is almost the same for early and late bookings.}\hfill
\subfig{0.26}{fig17.png}{4.5cm}{Test 5: the Business share is only slightly higher for last-minute bookings.}
\caption{Plots for the hypothesis tests. Each panel has its one-line takeaway.}
\label{fig:tests}
\end{figure*}

%------------------------------------------------------------ Test 2
\needspace{22\baselineskip}\subsection*{Test 2 (P2): are late bookings more expensive?}

\tpart{a}{Hypotheses, written before running the test}
{\small
\begin{tabularx}{\linewidth}{@{}K{3.4cm} Y@{}}
\toprule
$H_0$ & $\mu_{\text{late}}-\mu_{\text{early}}=0$ \\
$H_1$ & $\mu_{\text{late}}-\mu_{\text{early}}\ne0$ \\
Significance level & $\alpha=0.05$ \\
Side & two-sided \\
Test & two-sample z-test \\
Smallest difference that matters & \rs500 \\
\bottomrule
\end{tabularx}}

\egblock{Two-sided, because a late \emph{discount} (as on the last day) was also plausible.}

\tpart{b}{Other tests we considered}
\egblock{Both groups are very large (112,403 and 185,517), so the sample SDs estimate $\sigma$
almost exactly and the z statistic is valid without assuming normality.}
{\footnotesize
\begin{tabularx}{\linewidth}{@{}L{2.4cm} Y@{}}
\toprule
\textbf{Also considered} & \textbf{Why we did not use it} \\
\midrule
Welch t-test & identical at this size: $t=46.09$ (239,268); adds nothing \\
Mann--Whitney U & checks the result is not driven by outliers: $p<10^{-300}$, $P(\text{late}>\text{early})=0.65$, medians \rs11,511 vs.\ \rs6,067 \\
Paired t-test & nothing to pair: listings are not matched across windows \\
One-way ANOVA over 49 days & answers a vaguer question (``some day differs''); Test 4 models the trend instead \\
z-test within each class & run as a check, since class dominates price (result below) \\
\bottomrule
\end{tabularx}}

\tpart{c}{Result}
{\small
\begin{tabularx}{\linewidth}{@{}K{3.4cm} Y@{}}
\toprule
Estimate & \rs3,924 (\rs23,327 vs.\ \rs19,402), +20.2\% \\
95\% confidence interval & \rs3,757 to \rs4,091 \\
Test statistic (d.f.) & $z=46.09$ (none for a z-test) \\
p-value & $<10^{-300}$ \\
Decision & reject $H_0$ \\
\bottomrule
\end{tabularx}}

\egblock{Late bookings cost about \rs3.9k more, far above \rs500 (Fig.~\ref{fig:tests}b). The
premium is \rs4,069 [4,036, 4,102] (+81\%) within Economy and \rs1,846 [1,673, 2,018] (+3.6\%)
within Business. Narrowing ``late'' to the last week widens the gap to \rs6,912 ($z=49.4$).}

%------------------------------------------------------------ Test 3
\needspace{22\baselineskip}\subsection*{Test 3 (P2): are late-booking fares more volatile?}

\tpart{a}{Hypotheses, written before running the test}
{\small
\begin{tabularx}{\linewidth}{@{}K{3.4cm} Y@{}}
\toprule
$H_0$ & $\sigma^2_{\text{late}}/\sigma^2_{\text{early}}\le1$ \\
$H_1$ & $\sigma^2_{\text{late}}/\sigma^2_{\text{early}}>1$ \\
Significance level & $\alpha=0.05$ \\
Side & one-sided (right) \\
Test & F-test for two variances \\
Smallest difference that matters & SD 10\% higher (variance ratio 1.21) \\
\bottomrule
\end{tabularx}}

\tpart{b}{Other tests we considered}
\egblock{P2 asks whether the late variance is \emph{larger}; the F ratio is exactly that, and it
allows a one-sided alternative and a CI for the ratio. It assumes normality, which fails
(Fig.~\ref{fig:eda}h), so we compared it with robust tests.}
{\footnotesize
\begin{tabularx}{\linewidth}{@{}L{2.4cm} Y@{}}
\toprule
\textbf{Also considered} & \textbf{Why we did not use it} \\
\midrule
Bartlett & also assumes normality: $\chi^2=18.6$, $p=2\times10^{-5}$ \\
Levene (mean) & robust but two-sided only: $W=134.2$ (1, 297,918), $p=5\times10^{-31}$ \\
Brown--Forsythe (median) & most robust, two-sided only: $W=169.9$, $p=8\times10^{-39}$; mean $|x-\text{median}|$ is \rs16,328 (late) vs.\ \rs15,309 (early) \\
F-test within each class & run as a check, since the class gap dominates the pooled spread \\
\bottomrule
\end{tabularx}}

\tpart{c}{Result}
{\small
\begin{tabularx}{\linewidth}{@{}K{3.4cm} Y@{}}
\toprule
Estimate & variance ratio 0.977 (SD \rs22,430 vs.\ \rs22,690) \\
95\% confidence interval & 0.967 to 0.988 \\
Test statistic (d.f.) & $F=0.977$ (112,402; 185,516) \\
p-value & 1.000 \\
Decision & fail to reject $H_0$ \\
\bottomrule
\end{tabularx}}

\egblock{Pooled over both classes, late fares are not more variable: the interval lies below 1,
far from 1.21 (Fig.~\ref{fig:tests}f). The tests disagree because they measure different things,
and pooling hides real volatility. Within Economy, late fares have more than twice the SD
(\rs4,425 vs.\ \rs2,030; $F=4.75$ [4.69, 4.81]); within Business they are about 10\% more spread
($F=1.21$ [1.19, 1.23]); both $p<10^{-90}$. P2's answer is \emph{yes within each class, no when
pooled}.}

%------------------------------------------------------------ Test 4
\needspace{22\baselineskip}\subsection*{Test 4 (P1): does fare fall as booking moves earlier?}

\tpart{a}{Hypotheses, written before running the test}
{\small
\begin{tabularx}{\linewidth}{@{}K{3.4cm} Y@{}}
\toprule
$H_0$ & $\beta_1=0$ in $\text{price}=\beta_0+\beta_1\,\text{days\_left}+\varepsilon$ \\
$H_1$ & $\beta_1\ne0$ \\
Significance level & $\alpha=0.05$ \\
Side & two-sided \\
Test & t-test on the OLS slope \\
Smallest difference that matters & \rs100 per day (\rs1,000 for 10 days) \\
\bottomrule
\end{tabularx}}

\tpart{b}{Other tests we considered}
\egblock{P1 is a regression problem; the slope is the average saving per extra day of advance
booking, and its t-test is the standard inference for it.}
{\footnotesize
\begin{tabularx}{\linewidth}{@{}L{2.4cm} Y@{}}
\toprule
\textbf{Also considered} & \textbf{Why we did not use it} \\
\midrule
HC1 robust SE & residuals are heteroscedastic; robust SE 3.06 vs.\ 3.05, $t=-49.7$, same decision \\
Spearman, Kendall & test any monotone trend, not the size of the saving: $\rho=-0.266$, $\tau_b=-0.185$, both $p<10^{-300}$; much stronger than Pearson $r=-0.091$, so the trend is monotone but not linear \\
Regression with class & checks the slope is not a class artefact: $\hat\beta_1=-130.6$ (SE 1.03), $R^2=0.887$ \\
Piecewise slopes & test the sweet spot directly (result below) \\
\bottomrule
\end{tabularx}}

\tpart{c}{Result}
{\small
\begin{tabularx}{\linewidth}{@{}K{3.4cm} Y@{}}
\toprule
Estimate & $-$\rs152 per extra day booked ahead \\
95\% confidence interval & $-157.9$ to $-146.0$ \\
Test statistic (d.f.) & $t=-49.77$ (297,918) \\
p-value & $<10^{-300}$ \\
Decision & reject $H_0$ \\
\bottomrule
\end{tabularx}}

\egblock{The whole interval is beyond \rs100/day (Fig.~\ref{fig:tests}c). Pooled $R^2$ is only
0.008 because class dominates; within Economy it is 0.313. The piecewise fit locates the sweet
spot: within 20 days of departure each extra day booked ahead saves \rs464 [440, 488], but beyond
20 days only \rs25 [13, 38], under the \rs100 threshold. \textbf{Booking about three weeks ahead
captures nearly all the savings.}}

%------------------------------------------------------------ Test 5
\needspace{22\baselineskip}\subsection*{Test 5 (P2): does the class mix change for last-minute bookings?}

\tpart{a}{Hypotheses, written before running the test}
{\small
\begin{tabularx}{\linewidth}{@{}K{3.4cm} Y@{}}
\toprule
$H_0$ & booking timing (last-minute $\le7$ days / standard) is independent of class \\
$H_1$ & the two are associated \\
Significance level & $\alpha=0.05$ \\
Side & right tail of $\chi^2$ \\
Test & Pearson $\chi^2$ test of independence \\
Smallest difference that matters & 5 percentage points in Business share \\
\bottomrule
\end{tabularx}}

\egblock{This checks whether the premium in Tests 1--2 could simply be more Business tickets among
late bookings.}

{\footnotesize
\begin{tabularx}{\linewidth}{@{}L{2.6cm} Y Y Y@{}}
\toprule
 & \textbf{Business} & \textbf{Economy} & \textbf{Business share} \\
\midrule
Last-minute ($\le7$ d) & 10,422 & 21,211 & 32.95\% \\
Standard ($>7$ d) & 82,706 & 183,581 & 31.06\% \\
\bottomrule
\end{tabularx}}

\tpart{b}{Other tests we considered}
\egblock{Two categorical variables and all expected counts $\ge9{,}888$, so the $\chi^2$
approximation is excellent.}
{\footnotesize
\begin{tabularx}{\linewidth}{@{}L{2.4cm} Y@{}}
\toprule
\textbf{Also considered} & \textbf{Why we did not use it} \\
\midrule
Yates-corrected $\chi^2$ & correction matters only for small counts: $\chi^2=46.80$, same decision \\
G-test & likelihood-ratio version, equivalent at this size: $G=46.50$, $p=9\times10^{-12}$ \\
Two-proportion z-test & equivalent ($z=6.85$, $z^2=\chi^2$); we use it for the CI below \\
Fisher's exact test & built for small expected counts, so unnecessary: OR $=1.09$, $p=9\times10^{-12}$ \\
\bottomrule
\end{tabularx}}

\tpart{c}{Result}
{\small
\begin{tabularx}{\linewidth}{@{}K{3.4cm} Y@{}}
\toprule
Estimate & +1.89 pp Business share (Cramér's $V=0.013$) \\
95\% confidence interval & +1.34 to +2.43 pp \\
Test statistic (d.f.) & $\chi^2=46.89$ (1) \\
p-value & $7.5\times10^{-12}$ \\
Decision & reject $H_0$ \\
\bottomrule
\end{tabularx}}

\egblock{The association is real but does not matter: the whole interval is below 5 pp
(Fig.~\ref{fig:tests}g). So the late-booking premium in Tests 1--2 is not a class-mix effect.}


\minihead{What we conclude}
\egblock{Fares are flat beyond about 20 days and rise steeply inside that window, so about three
weeks ahead is the sweet spot (Test 4). Late bookings cost \rs3.9k more, +81\% in Economy
(Test 2), yet fares drop by \rs8.7k on the very last day (Test 1). Late fares are more volatile
within each class, which a pooled test misses (Test 3), and none of this is a class-mix effect
(Test 5).}

%-----------------------------------------------------------------------
\section{Limitations}

\begin{itemize}
  \item \textbf{Association, not cause.} Fare gaps reflect demand, competition and seat
  inventory, none of which is recorded; the tests cannot show that booking later \emph{causes}
  higher fares.
  \item \textbf{Listings are not fully independent.} The same flight appears on many search
  dates, so the effective sample is smaller than 297,920 and our p-values are optimistic. The
  effect sizes and intervals are less affected.
  \item \textbf{One 50-day window.} Seasonality, holidays, fuel surcharges and load factors are
  absent, so the 20-day sweet spot may differ in other seasons.
  \item \textbf{Dropped rows.} The 2,233 outlier rows (long layovers and extreme Business fares)
  are not covered by our conclusions.
  \item \textbf{Bimodal fares.} Business fares dominate the pooled variance and mask
  within-class effects (Tests 3--4); the P1 model must condition on class or use
  $\log(\text{price})$.
\end{itemize}

\end{document}
